% !PRINTED none
% Abschnitt-opening (part-title) page: no folio is printed on this page in the scan
% (section-opening pages in this book suppress the running head/page number).
% Content: "Sechzehnter Abschnitt." / "Kreistheilung." (both letter-spaced, centred),
% a horizontal divider rule, then "§. 167." and its letter-spaced title, followed
% immediately by the start of ordinary body text (no page break before the body begins).
\begin{center}
\gesperrt{Sixteenth Section.}

\gesperrt{Division of the Circle.}

\rule{3cm}{0.4pt}

\S.\,167.

\gesperrt{The cyclotomic periods and the period equations.}
\end{center}

The most important among the Abelian equations are those on which the determination of the roots of unity depends, whose elementary properties we have already come to know in the twelfth section. On account of the relation to the geometric problem of dividing the circumference of the circle into a definite number of equal parts, all the equations connected with this are also called the \gesperrt{cyclotomic equations}.

The field $\Omega$, which contains the quantities regarded as known, is here simply the \gesperrt{field of rational numbers}, which we shall call the field $R$.

We shall first occupy ourselves with the $n$th roots of unity under the assumption that $n$ is an odd prime (the prime 2 gives rise to no further questions, since the only second roots of unity, $\pm 1$, are contained in $R$).

Besides the rational $n$th root of unity 1 there exist $n - 1$ more primitive ones, which can be represented by the transcendental expressions
\[
e^{\frac{2\pi i}{n}}, \quad e^{\frac{4\pi i}{n}}, \quad e^{\frac{6\pi i}{n}} \dots e^{\frac{2(n-1)\pi i}{n}}
\]
If we denote one of them by $r$, then the whole system is also represented by the following list:
\weq{1}{r, \ r^2, \ r^3 \dots r^{n-1}}
%CONT>>
